% ============================================
% Johns Hopkins Letter Template
% Assistant Professor Cover Letter – University of Hawai‘i at Mānoa (SSRI/UHERO)
% ============================================

\documentclass[11pt,letterpaper]{letter}

\usepackage{graphicx}
\usepackage{eso-pic}
\usepackage{geometry}
\usepackage{microtype}
\usepackage[hidelinks]{hyperref}

% ----- Page Setup -----
\geometry{left=25mm,right=25mm,top=25mm,bottom=25mm}

% ----- Logo Placement (foreground, first page only) -----
\newcommand{\PutJHULogoFG}[1]{%
  \AddToShipoutPictureFG*{%
    \AtPageUpperLeft{%
      \hspace{10mm}%
      \raisebox{-38mm}{%
        \includegraphics[width=6cm]{#1}%
      }%
    }%
  }%
}

% ----- Sender Information -----
\signature{Decory Edwards}
\address{Department of Economics \\ Johns Hopkins University \\ Baltimore, MD 21218 \\ \texttt{dedwar65@jh.edu}}

\begin{document}
\PutJHULogoFG{Images/jhulogo.png}

\begin{letter}{Search Committee \\
Social Science Research Institute (SSRI) \\
University of Hawai‘i at Mānoa \\
Honolulu, HI 96822 \\ USA}

\opening{Dear Members of the Search Committee,}

I am writing to express my interest in the Assistant/Associate Professor position jointly housed in the Social Science Research Institute (SSRI) and the University of Hawai‘i Economic Research Organization (UHERO). I am a Ph.D. candidate in Economics at Johns Hopkins University, advised by Professors Christopher D. Carroll, Nicholas W. Papageorge, and Stelios Fourakis, and I expect to receive my degree in May 2026. My research fields are macroeconomics, wealth and income dynamics, and computational economics, with strong connections to empirical macroeconomics, regional economics, and applied econometric modeling.

My research agenda focuses on understanding the determinants of wealth inequality and the mechanisms through which heterogeneity in returns to wealth shapes aggregate outcomes. In my job-market paper, I estimate the degree of heterogeneity in individual portfolio returns using micro-data from the Survey of Consumer Finances and simulate a structural model in which households face idiosyncratic return risk. I show that allowing for persistent heterogeneity in returns substantially improves the model’s ability to match the empirical distribution of wealth, offering a tractable way to connect micro-level return dispersion to macroeconomic inequality.

In a second project, I use the Health and Retirement Study to examine the relationship between interpersonal trust and portfolio performance. Building on the literature that finds a hump-shaped relationship between trust and income, I show that a similar relationship holds for individual returns to wealth measured in the HRS data, highlighting a behavioral channel through which social attitudes shape saving and investment outcomes. This work also provides new estimates of the persistent component of returns to net worth, which motivate the calibration of heterogeneity in my structural model.

The mission of UHERO and SSRI—to provide rigorous, policy-relevant empirical analysis on issues central to Hawai‘i’s economy, including business cycles, regional dynamics, tourism, housing, sustainability, and long-run economic resilience—resonates deeply with my research approach. My work integrates microdata, structural modeling, and econometric tools to address questions at the intersection of macroeconomic fluctuations, inequality, investment behavior, and household financial resilience. I am eager to contribute to UHERO’s forecasting efforts, applied econometric modeling, and policy analysis, and to collaborate with researchers across the College of Social Sciences on projects related to regional macroeconomics, tourism economics, climate and environmental shocks, and data-driven policy evaluation. The opportunity to work in an applied, interdisciplinary environment that bridges academic research with real-time policy relevance is exceptionally well aligned with my training and goals.

\textbf{Teaching.} My teaching experience centers on undergraduate macroeconomics and an upper-level macro-policy seminar, where I have led weekly discussion sections and held regular office hours for classes ranging from 16 to 40 students. My approach is student-centered: I aim to help students “convince themselves” that they have moved from confusion to understanding by holding structured office-hour coaching that builds trust and confidence. At UH Mānoa, I am prepared to teach \textit{Principles of Macroeconomics}, \textit{Intermediate Macroeconomic Theory}, \textit{Money and Banking}, \textit{Econometrics}, and \textit{Quantitative Methods for Economics}. I would also welcome developing electives such as \textit{Computational Macroeconomics and Policy Modeling} or \textit{Economics of Inequality and Wealth Dynamics}, incorporating reproducible workflows and real Hawai‘i-focused datasets.

I am committed to fostering an inclusive learning environment consistent with the University of Hawai‘i’s mission and Hawai‘i’s multicultural community. In my teaching and mentorship, I integrate diverse data sources, accessible computational tools, and multiple modes of assessment to ensure that students from varied backgrounds can participate fully in economics. I also aim to expand undergraduate research opportunities—particularly for first-generation and underrepresented students—by incorporating coding workshops, replication exercises, and structured research mentoring.

I would be honored to join SSRI and UHERO and contribute to their research, teaching, and policy missions. Thank you for your consideration; I welcome the opportunity to discuss my fit with your program at your convenience.

\closing{Sincerely,}

\end{letter}
\end{document}
